\chapter{Conclusion and Future Work}
\label{chap:conclusion}

This thesis investigated the integration of Federated Dynamic Averaging (FDA) with Parameter-Efficient Low-Rank Adaptation (LoRA) in federated learning. By formulating and analyzing the FDA-RoLoRA Dual Track protocol, this work examined how variance-driven dynamic synchronization can be coupled with alternating low-rank parameter optimization to jointly address spatial parameter overhead and rigid temporal synchronization schedules under statistical data heterogeneity.

\section{Synthesis of Results and Methodological Insights}
\label{sec:synthesis_results}

The development and evaluation of the FDA-RoLoRA Dual Track framework yield several key theoretical and empirical insights:

\paragraph{Mathematical Consistency Through Alternating Subspace Isolation.}
A central insight of this work is that dynamic averaging cannot be naively paired with standard federated LoRA due to the bilinear aggregation inconsistency ($\frac{1}{N}\sum A_i B_i \neq (\frac{1}{N}\sum A_i)(\frac{1}{N}\sum B_i)$). Bilinear aggregation errors distort client drift trajectories, causing geometric variance monitoring to misestimate true model divergence. By adopting RoLoRA's alternating coordinate scheme, the active parameter space is restricted to a single factor ($A$ or $B$) during each communication round. This guarantees strict linearity in local updates, providing the exact mathematical property required for stable, decoupled subspace variance tracking without cross-matrix distortion.

\paragraph{Subspace Variance Dynamics and Temporal Smoothing.}
Evaluating variance directly within low-rank coordinate tensors eliminates the memory and computational overhead of full-model gradient approximations. However, low-rank parameter trajectories exhibit significantly higher gradient volatility than full-model fine-tuning. Unsmoothed threshold prediction based on instantaneous variance often leads to aggressive round termination or boundary clipping at $e_{\max}$. Incorporating temporal smoothing—specifically Simple Moving Average (SMA) and Exponential Moving Average (EMA)—proved essential for stabilizing threshold evolution across rounds, balancing local epoch utilization, and preventing premature synchronization.

\paragraph{Empirical Trade-Offs on Synthetic Benchmarks.}
Extensive empirical evaluation across 66 runs on a controlled synthetic benchmark under extreme non-IID data partitioning substantiated the core behavior of the protocol:
\begin{itemize}
    \item \textbf{Communication Efficiency and Reliability}: Dynamic averaging variants demonstrated faster convergence to target accuracy milestones (achieving up to a $2.20\times$ conditional round reduction on non-linear partitions) and exhibited higher convergence reliability than fixed-epoch baselines.
    \item \textbf{Compute versus Communication Trade-Off}: Dynamic averaging expended $1.7\times$ to $2.2\times$ more local gradient steps than fixed-schedule alternatives. This confirms that dynamic synchronization provides the highest practical benefit in federated settings where communication latency, round orchestration, or network bandwidth dominate overall training time, whereas fixed-budget methods remain preferable when edge computation is strictly constrained.
\end{itemize}

\paragraph{Scope and Status of Foundation Model Evaluations.}
While the multi-process software architecture and model injection pipelines were fully implemented for large foundation models (RoBERTa-Large and Llama-2-7B on GLUE classification tasks), execution of these large-scale experiments remains in progress on cluster resources at the time of writing. Consequently, the synthetic benchmark results serve as the primary empirical validation of the dual-track synchronization mechanism and smoothing dynamics.

\section{Limitations}
\label{sec:limitations}

While this work demonstrates the feasibility and properties of combining dynamic averaging with alternating low-rank adaptation, several operational and methodological limitations warrant discussion:

First, from an empirical standpoint, the primary completed experimental results are derived from a controlled synthetic classification benchmark, with large-scale foundation model evaluations on GLUE tasks currently queued. While the synthetic setup captured extreme label skew and multi-layer non-linear dynamics, empirical validation across billion-parameter language models is needed to confirm scaling behaviors.

Second, the lack of access to the original codebase from the RoLoRA paper mandated a complete custom reimplementation within our framework. Furthermore, experimental evaluations were conducted entirely within a simulated multi-process federated environment on institutional GPU resources ($2\times$ NVIDIA A10 GPUs), leaving real-world distributed network phenomena—such as physical link latency, bandwidth fluctuations, and stochastic client dropouts—unmodeled.

Third, regarding variance tracking, while exact variance computation directly on compact LoRA tensors proved efficient for low adapter ranks ($r=8$), scaling to significantly higher ranks or wider layer coverage may eventually require approximate sketching techniques (such as the AMS sketch). Finally, the alternating optimization protocol relies on a rigid alternating schedule (e.g., $B$-first alternation), leaving the exploration of data-adaptive or variance-driven matrix toggling criteria for future investigation.

\section{Future Work}
\label{sec:future_work}

Building upon the methodology and findings established in this work, several promising avenues for future research emerge:

\paragraph{Completion of Foundation Model Benchmarking.}
The immediate priority is completing the execution of the large-scale evaluation suite on RoBERTa-Large ($355\text{M}$) and Llama-2-7B ($7\text{B}$) across GLUE benchmark tasks (SST-2, MNLI) on the Polytechnix cluster. This will provide definitive empirical data on downstream task accuracy, total communication rounds, and byte transmission under realistic language modeling workloads.

\paragraph{Dynamic Averaging Across Other Federated LoRA Variants.}
Future work should explore extending dynamic variance monitoring to alternative federated parameter-efficient techniques. Promising directions include investigating \textbf{FDA + SLoRA} \cite{babakniya2023slora} for sparse low-rank coordinate exchanges, \textbf{FDA + FFA-LoRA} \cite{sun2024improving} with permanently frozen down-projections, and analyzing whether dynamic variance monitoring can help diagnose and mitigate client drift in standard simultaneous \textbf{FDA-FedAvg-LoRA}. Furthermore, developing adaptive rank selection—where the rank $r$ dynamically expands or contracts based on observed subspace variance—represents a natural extension.

\paragraph{Optimization and System Extensions.}
From an optimization perspective, exploring asymmetric learning rate schedules tailored independently for down-projection ($A$) and up-projection ($B$) matrices offers potential convergence benefits. Concurrently, developing formal convergence proofs for alternating dynamic averaging under non-convex loss surfaces constitutes an important theoretical frontier. System-level integrations with low-bit quantization (such as QLoRA \cite{dettmers2023qlora}) and deployment on physical edge testbeds will further assess operational feasibility.

\paragraph{Integration with Federated Learning Frameworks.}
To encourage community adoption and standardized benchmarking, integrating the FDA-RoLoRA Dual Track protocol into established open-source federated benchmarking platforms—such as \textbf{FederatedScope-LLM} \cite{kuang2023federatedscope}—will facilitate broader evaluation across diverse model architectures and distributed settings.

\paragraph{Concluding Remarks.}
In conclusion, this work demonstrates that federated low-rank adaptation can effectively incorporate dynamic, variance-driven synchronization. By establishing that alternating coordinate optimization provides a mathematically consistent substrate for geometric variance monitoring, this thesis provides a principled and communication-efficient approach for federated fine-tuning.

